\documentclass[11pt]{article}

\usepackage[a4paper,margin=27mm]{geometry}
\usepackage[T1]{fontenc}
\usepackage{lmodern}
\usepackage{amsmath,amssymb,amsthm,mathtools}
\usepackage{microtype}
\usepackage{xcolor}
\usepackage[colorlinks=true,linkcolor=blue!55!black,
  urlcolor=blue!55!black]{hyperref}

\newtheorem{theorem}{Theorem}[section]
\newtheorem{lemma}[theorem]{Lemma}
\newtheorem{proposition}[theorem]{Proposition}
\newtheorem{corollary}[theorem]{Corollary}
\theoremstyle{definition}
\newtheorem{definition}[theorem]{Definition}
\newtheorem{construction}[theorem]{Construction}
\theoremstyle{remark}
\newtheorem{remark}[theorem]{Remark}

\newcommand{\R}{\mathbb R}
\newcommand{\one}{\mathbf 1}
\newcommand{\floor}[1]{\left\lfloor #1\right\rfloor}

\title{The no-three-collinear variant V3:\\[1mm]a compact, search-free proof}
\author{Companion note}
\date{5 August 2026}

\begin{document}
\maketitle

\begin{abstract}
For \(n\geq4\), let \(f_{\mathrm{V3}}(n)\) be the least number of proper
Euclidean circles determined by \(n\) distinct real planar points, no three
collinear and not all concyclic.  We prove
\[
 f_{\mathrm{V3}}(8)=20,
 \qquad
 f_{\mathrm{V3}}(n)=1+\binom{n-1}{2}\quad(n\neq8).
\]
The infinite range follows from a summed inverted Melchior inequality and
two explicit lower bounds.  The only equality obstructions left in the
small range are proved here: a seven-point Fano wall and a ten-point
six-circle power wall.  No configuration, support, or orbit catalogue and
no SAT, ILP, or finite-range search is used.
\end{abstract}

\section{Statement and geometric conventions}

\begin{definition}[V3]
A \emph{V3 configuration} is a set \(P\subset\R^2\) of \(n\geq4\) distinct
points such that no three points of \(P\) are collinear and \(P\) is not
contained in one proper circle.  A circle is counted when it is a proper
Euclidean circle containing at least three points of \(P\), and it is
counted only once, independently of the size of its full trace on \(P\).
We denote the count by \(C(P)\) and its minimum by \(f_{\mathrm{V3}}(n)\).
\end{definition}

Put
\begin{equation}\label{eq:G}
 G(n):=1+\binom{n-1}{2}.
\end{equation}

We use the following standard facts, always over the real plane.  A
noncollinear triple has a unique circumcircle; two distinct proper circles,
or a proper circle and a line, meet in at most two points.  Inversion about
\(o\) interchanges the lines avoiding \(o\) with the proper circles through
\(o\), sends circles avoiding \(o\) to circles avoiding \(o\), and preserves
all incidences away from \(o\).  For a proper circle and two secants through
a point, directed power of a point equates the two products of directed
intercepts.  Directed Menelaus applies to a nondegenerate real affine
triangle and three collinear points on its extended sides; none of the six
ratios used below has a zero denominator.  The one global incidence input,
Melchior's inequality, is derived in \S\ref{sec:universal} for every finite
noncollinear real projective point set.  Thus all classical inputs are used
on their stated real, finite, nondegenerate ranges.

\begin{theorem}\label{thm:v3}
For every integer \(n\geq4\),
\[
 \boxed{
 f_{\mathrm{V3}}(n)=
 \begin{cases}
 20,&n=8,\\[1mm]
 G(n),&n\neq8.
 \end{cases}}
\]
In particular,
\[
 \bigl(f_{\mathrm{V3}}(4),\ldots,f_{\mathrm{V3}}(10)\bigr)
   =(4,7,11,16,20,29,37).
\]
\end{theorem}

\section{The two upper constructions}

\begin{construction}[A circle and one outsider]\label{con:near-circle}
Choose \(n-1\) distinct points on a proper circle \(\Gamma\), and choose
\(x\notin\Gamma\) outside every line through two of them.  The forbidden set
is a finite union of proper curves, so such an \(x\) exists.  The resulting
set has no collinear triple.  It is not concyclic, since any circle through
the \(n-1\geq3\) selected points of \(\Gamma\) equals \(\Gamma\), while
\(x\notin\Gamma\).  Thus it is V3.  Its triples on \(\Gamma\) give
\(\Gamma\), and every other triple
consists of \(x\) and a pair on \(\Gamma\).  Different pairs give different
circles: a coincidence would produce a circle through \(x\) sharing at least
three points with \(\Gamma\).  Hence
\begin{equation}\label{eq:upper-G}
 f_{\mathrm{V3}}(n)\leq1+\binom{n-1}{2}=G(n).
\end{equation}
\end{construction}

The exceptional construction is just as explicit, but it is clearer before
inversion, where lines and circles are temporarily treated together as
\emph{generalized circles}.

\begin{construction}[The eight-point cross]\label{con:cross}
Fix \(0<a<b\), let \(S=\{-b,-a,a,b\}\), and set
\[
 Q=(S\times\{0\})\cup(\{0\}\times S).
\]
The two axes are four-point generalized circles.  If \(x_1\neq x_2\) and
\(y_1\neq y_2\), then
\begin{equation}\label{eq:product-criterion}
 (x_1,0),(x_2,0),(0,y_1),(0,y_2)
 \text{ are concyclic}
 \quad\Longleftrightarrow\quad x_1x_2=y_1y_2.
\end{equation}
Indeed, on restricting
\(X^2+Y^2+DX+EY+F=0\) to the two axes, Vieta's formula says that both
root pairs have product \(F\); conversely, equality of the products fixes
\(D,E,F\) and gives a proper circle.

The six unordered pairs from \(S\) have four distinct product fibres,
\[
 \begin{array}{c|cccc}
 \text{product}&ab&-ab&-b^2&-a^2\\ \hline
 \text{multiplicity}&2&2&1&1.
 \end{array}
\]
Consequently \eqref{eq:product-criterion} gives
\(2^2+2^2+1^2+1^2=10\) four-point \(2+2\) circles.  There are \(48\) mixed
triples.  Those ten circles own \(10\binom43=40\) of them.  Each of the
remaining eight triples has its own three-point circle: if two had the same
carrier, that carrier could meet each axis at most twice, so their union
would be a four-point \(2+2\) support and hence one of the ten circles
already counted.  Thus the generalized blocks are
\[
 2\text{ axes},\qquad10\text{ four-point circles},\qquad
 8\text{ three-point circles},
\]
and both the block and triple counts are exact:
\begin{equation}\label{eq:cross-count}
 2+10+8=20,
 \qquad
 2\binom43+10\binom43+8=\binom83.
\end{equation}

Let \(\mathcal U\) be the union of these twenty lines and circles.  A finite
union of such curves does not cover the plane, so choose \(o\notin\mathcal U\)
and invert about \(o\).  Since \(Q\subset\mathcal U\), the centre is not a
selected point.  Every one of the twenty generalized blocks becomes a
proper circle with the same selected trace.  Conversely, a circle through
an image triple pulls back to the unique generalized carrier of the
original triple, so there are no additional blocks.

No three image points are collinear: the inverse of such a line would be a
generalized circle containing \(o\) and an original triple, contrary to
\(o\notin\mathcal U\).  The image is not concyclic either, since its inverse
would put all of \(Q\) on one generalized circle, whereas a proper circle
meets either axis at most twice and no line contains all eight points.
Therefore
\begin{equation}\label{eq:upper-eight}
 f_{\mathrm{V3}}(8)\leq20.
\end{equation}
\end{construction}

The product fibres, not a list of supports, explain the exception: after
inversion the two rich axes and ten \(2+2\) collisions compress all
\(\binom83\) triples into only twenty proper circles.

\section{Universal lower bounds}\label{sec:universal}

Fix a V3 configuration.  Let \(c_s\) be the number of circles whose full
selected trace has size \(s\), let \(C=\sum_s c_s\), and let
\(m=\max\{s:c_s>0\}\).  Unique circumcircles partition the selected triples,
and nonconcyclicity gives \(3\leq m\leq n-1\).  Hence
\begin{equation}\label{eq:T}
 \boxed{\mathsf T:\quad
 \sum_{s=3}^{m}\binom{s}{3}c_s=\binom n3},
 \qquad C=\sum_{s=3}^{m}c_s.
\end{equation}

\subsection{Summed inverted Melchior}

For a finite noncollinear real projective point set, let \(t_r\) count its
spanned lines containing exactly \(r\) selected points.  Dualize to an
essential arrangement of real projective lines.  If \(V,E,F\) denote its
numbers of vertices, edges, and faces, then
\[
 V=\sum_{r\geq2}t_r,\qquad E=\sum_{r\geq2}r t_r,
 \qquad V-E+F=1.
\]
Every face has at least three sides, so \(3F\leq2E\).  Therefore
\(E\leq3V-3\), which is exactly
\begin{equation}\label{eq:melchior}
 t_2\geq3+\sum_{r\geq4}(r-3)t_r.
\end{equation}

Invert now at a selected point \(p\) and remove \(p\).  The \(n-1\) images
are noncollinear: a common image line through the centre would give a
collinear original triple, while a common image line avoiding the centre
would put all original points on one circle through \(p\).  Also, every
spanned image line avoids the centre, again by the no-three-collinear
hypothesis.  Thus an \(s\)-point circle through \(p\) becomes an
\((s-1)\)-point line, with no exceptional line term.

Sum \eqref{eq:melchior} over all \(n\) pivots.  A three-circle contributes
to \(t_2\) at each of its three points; a four-circle has coefficient zero;
and an \(s\)-circle, \(s\geq5\), contributes \(s-4\) at each of its \(s\)
points.  We obtain
\begin{equation}\label{eq:M}
 \boxed{\mathsf M:\quad
 3c_3-\sum_{s=5}^{m}s(s-4)c_s\geq3n.}
\end{equation}

\subsection{The bounded-support functional}

Assume \(m\geq4\), and put
\[
 A=\binom m3,\qquad B=m(m-4),\qquad D=3A+B,
 \qquad
 \lambda_m=\frac{B+3}{D},\quad \mu_m=\frac{A-1}{D}.
\]
Here \(D=m(m-3)(m+2)/2>0\), and both multipliers are positive.  Add
\(\lambda_m\mathsf T\) and \(\mu_m\mathsf M\).  Directly from the
definitions,
\[
 \lambda_m+3\mu_m=1,\qquad
 \lambda_m\binom m3-\mu_m m(m-4)=1,
\]
so the coefficients of \(c_3\) and \(c_m\) are one.  For every integer
\(4\leq s\leq m\), the exact slack below one is
\begin{equation}\label{eq:slack}
 1-\left(\lambda_m\binom s3-\mu_m s(s-4)\right)
 =\frac{(m-s)(s-3)(ms-m-s-2)}{3m(m+2)}\geq0,
\end{equation}
because \(ms-m-s-2=(m-1)(s-1)-3\geq6\).  The weighted left side is
therefore at most \(C\), whereas its right side is a lower bound.  Hence
\begin{equation}\label{eq:L}
 C\geq L(n,m):=\lambda_m\binom n3+3n\mu_m
 =\frac{n(3m^2+mn^2-3mn+2m-n^2+3n+4)}{3m(m+2)}.
\end{equation}
Direct subtraction gives the useful strict monotonicity
\begin{equation}\label{eq:L-monotone}
 L(n,m)-L(n,m+1)
 =\frac{n(n-4)(n+1)(m^2-m-3)}
 {3m(m+1)(m+2)(m+3)}>0
\end{equation}
for \(n\geq5\) and \(m\geq4\).

\subsection{The rich-circle pencil}

Suppose a circle \(\Gamma\) contains \(m\) selected points, with
\(k=n-m\) outsiders.  For each outsider \(x\), the circles through \(x\)
and a pair on \(\Gamma\) form a fan of exactly \(\binom m2\) circles.  A
circle common to the fans of \(x\) and \(y\) uses the same pair on
\(\Gamma\).  The pairs arising from two distinct common circles are
disjoint, since otherwise the circles would share \(x,y\) and a point of
\(\Gamma\).  Thus two fans meet in at most \(\floor{m/2}\) circles.
Two-term Bonferroni, together with \(\Gamma\), yields
\begin{equation}\label{eq:P}
 C\geq P(n,m):=
 1+(n-m)\binom m2-\binom{n-m}{2}\floor{\frac m2}.
\end{equation}

For later use suppose \(m>\floor{n/2}\), so \(1\leq k<m\).  If \(m=2r\),
then \(k\leq2r-1\) and
\begin{align}
 2(P(n,2r)-G(n))
  &=(k-1)\bigl(4r^2-6r+2-k(r+1)\bigr),\label{eq:P-even}\\
 4r^2-6r+2-k(r+1)
  &\geq(2r-1)(r-3).\nonumber
\end{align}
If \(m=2r+1\), then \(k\leq2r\) and
\begin{align}
 2(P(n,2r+1)-G(n))
  &=(k-1)\bigl(4r^2-2r-k(r+1)\bigr),\label{eq:P-odd}\\
 4r^2-2r-k(r+1)&\geq2r(r-2).\nonumber
\end{align}
For \(n\geq11\), every relevant even or odd rich case has \(r\geq3\).
Consequently
\begin{equation}\label{eq:P-rich}
 P(n,m)\geq G(n)
 \qquad(n\geq11, m>\floor{n/2}).
\end{equation}

\section{The whole infinite range}

Let \(n\geq11\).  If \(m=3\), then \eqref{eq:T} gives
\(C=\binom n3>G(n)\).  If \(m>\floor{n/2}\), use
\eqref{eq:P-rich}.  In the remaining range
\(4\leq m\leq\floor{n/2}\), monotonicity gives
\[
 C\geq L(n,m)\geq L\bigl(n,\floor{n/2}\bigr).
\]

For \(n=2r\), \(r\geq6\), exact subtraction yields
\begin{equation}\label{eq:large-even}
 L(2r,r)-G(2r)
 =\frac{(r-2)(2r^2-13r+2)}{3(r+2)}.
\end{equation}
At \(r=6\) this is \(-2/3\), and integrality of \(C\) and \(G(12)\) rounds
\(C\geq G(12)-2/3\) up to \(C\geq G(12)\).  At \(r=7\) the quadratic factor
is \(9\), and thereafter its forward difference \(4r-11\) is positive.

For \(n=2r+1\), \(r\geq5\), one obtains
\begin{equation}\label{eq:large-odd}
 L(2r+1,r)-G(2r+1)
 =\frac{(2r-3)(r^3-4r^2-4r-2)}{3r(r+2)}.
\end{equation}
The cubic factor is \(3\) at \(r=5\), and its forward difference
\(3r^2-5r-7\) is positive for \(r\geq5\).  The three ranges for \(m\) are
disjoint and exhaustive, so
\begin{equation}\label{eq:large-result}
 C\geq G(n)\qquad(n\geq11).
\end{equation}

\section{The two small real walls}

\subsection{Four-point circles and the Fano wall}

When \(m=4\), \eqref{eq:T} and \eqref{eq:M} reduce to
\begin{equation}\label{eq:m4}
 \binom n3=c_3+4c_4,\qquad c_3\geq n.
\end{equation}
Let \(\rho(n)\) be the least integer at least \(n\) congruent to
\(\binom n3\pmod4\).  Then
\begin{equation}\label{eq:m4-bound}
 C=\frac{\binom n3+3c_3}{4}
 \geq\frac{\binom n3+3\rho(n)}4.
\end{equation}
For \(n=5,\ldots,10\), the values of \(\rho(n)\) are
\[
 (6,8,7,8,12,12),
\]
and \eqref{eq:m4-bound} gives respectively
\begin{equation}\label{eq:m4-values}
 (7,11,14,20,30,39).
\end{equation}
Only the \(n=7\) equality lies below the eventual target.

\begin{lemma}[The planar Fano wall]\label{lem:fano}
Seven distinct real points cannot carry seven proper four-point circles
whose three-point complements form a Steiner triple system \(\mathrm{STS}(7)\).
\end{lemma}

\begin{proof}
First normalize the Steiner system without a classification.  Name one
triple \(123\).  The other triples through \(1\) pair the remaining four
points, so relabel them as \(145,167\).  The triples through \(2\) must use
the other pairings; after swapping \(6,7\) if needed, take \(246,257\).
The two remaining unused pairs among \(4,5,6,7\) are \(47,56\), forcing
\begin{equation}\label{eq:fano-lines}
 123,\quad145,\quad167,\quad246,\quad257,\quad347,\quad356.
\end{equation}

Invert at point \(1\).  The circles complementary to the last four triples
in \eqref{eq:fano-lines} become lines with full selected traces
\[
 357,\qquad346,\qquad256,\qquad247.
\]
The first three have the three distinct pairwise intersections \(3,5,6\);
they are therefore the sides of a nondegenerate triangle.  Set
\[
 A=3,\quad B=5,\quad C=6,\qquad
 X=7\in AB,\quad Y=4\in AC,\quad Z=2\in BC.
\]
The fourth line is the transversal \(XYZ\).  The three circles not through
the inversion centre have traces
\begin{equation}\label{eq:fano-circles}
 ABYZ,\qquad ACXZ,\qquad BCXY.
\end{equation}
All vertices and marked points are distinct because the displayed traces
are full traces of distinct carriers.

Use directed affine parameters
\[
 AB:(A,B,X)=(0,1,x),\quad
 AC:(A,C,Y)=(0,1,y),\quad
 BC:(B,C,Z)=(0,1,z),
\]
and put \(a^2=|BC|^2\), \(b^2=|CA|^2\), \(c^2=|AB|^2\).
Directed power at \(C,B,A\), respectively, for the circles in
\eqref{eq:fano-circles} gives
\[
 b^2(1-y)=a^2(1-z),\qquad
 c^2(1-x)=a^2z,\qquad
 c^2x=b^2y.
\]
The six factors \(x,y,z,1-x,1-y,1-z\) are nonzero.  Multiplying the three
resulting ratios gives
\begin{equation}\label{eq:fano-plus}
 \frac{x}{1-x}\frac{z}{1-z}\frac{1-y}{y}=+1.
\end{equation}
But directed Menelaus for the transversal \(XYZ\) in triangle \(ABC\)
gives the same product equal to \(-1\), a contradiction.
\end{proof}

Return to \(n=7,m=4\).  Congruence in \eqref{eq:m4} says either
\(c_3=7\) or \(c_3\geq11\).  In the former case \(c_4=7\).  If \(A_i\) are
the seven circle traces and \(B_i=P\setminus A_i\), then for \(i\neq j\)
\[
 |B_i\cap B_j|=|A_i\cap A_j|-1\leq1,
\]
because two four-subsets of a seven-set meet and two distinct circles meet
in at most two selected points.  Thus the seven triples \(B_i\) contain
\(7\binom32=21\) distinct pairs, exactly all pairs of \(P\); they form an
\(\mathrm{STS}(7)\), contrary to Lemma~\ref{lem:fano}.  Hence
\(c_3\geq11\), and \eqref{eq:m4-bound} improves to
\begin{equation}\label{eq:fano-jump}
 C\geq17\qquad(n=7,m=4).
\end{equation}

\subsection{Six five-point circles on ten points}

\begin{lemma}[The six-five wall]\label{lem:six-five}
Ten distinct real points cannot support six distinct proper five-point
circles if every two of their full traces meet in at most two points.
\end{lemma}

\begin{proof}
Let the traces be \(B_1,\ldots,B_q\) in a ten-set \(X\), and put
\(w_i=\one_{B_i}-\tfrac12\one_X\).  Then
\[
 \|w_i\|^2=\frac52,\qquad
 \langle w_i,w_j\rangle=|B_i\cap B_j|-\frac52\leq-\frac12.
\]
Therefore
\begin{equation}\label{eq:energy}
 0\leq\left\|\sum_{i=1}^q w_i\right\|^2
 \leq\frac{5q}{2}-\binom q2=\frac{q(6-q)}2.
\end{equation}
In particular \(q\leq6\).  If \(q=6\), equality throughout
\eqref{eq:energy} forces
\begin{equation}\label{eq:design-equality}
 |B_i\cap B_j|=2\quad(i\neq j),
 \qquad d(x)=3\quad(x\in X),
\end{equation}
where the second conclusion follows coordinatewise from
\(\sum_iw_i=0\).

We now derive, rather than assume, the equality design.  Delete \(B_6\).
Its five points have degree two among \(B_1,\ldots,B_5\), and the other five
have degree three.  Record the memberships of the first five as edges of a
multigraph \(E\) on block labels \(1,\ldots,5\); record the two-label
complements of the memberships of the latter five as edges of a multigraph
\(H\).  Each label has degree two in \(E\), because
\(|B_i\cap B_6|=2\).  It also has degree two in \(H\): among the five
points outside \(B_6\), label \(i\) occurs three times and is omitted twice.
Thus both graphs are \(2\)-regular.

The graph \(E\) has no doubled edge.  Otherwise two degree-two points would
exhaust \(B_i\cap B_j\), so no degree-three membership could contain both
\(i,j\).  Yet \(i,j\) occur six times in total among the five degree-three
memberships, impossible if each membership contains at most one of them.
The graph \(H\) has no doubled edge either.  If \(ij\) were doubled, the
degrees of \(i,j\) would be exhausted and the other three vertices would
form a triangle.  Its three edges are disjoint from \(ij\), so their
complementary memberships would put at least three points in
\(B_i\cap B_j\), contradicting \eqref{eq:design-equality}.  Hence \(E,H\)
are simple \(2\)-regular graphs on five vertices, and thus both are
five-cycles.

For a label pair \(ij\), the number of \(H\)-edges disjoint from it is
\(1+\one_{\{ij\in H\}}\).  Counting \(B_i\cap B_j\) now gives
\[
 2=\one_{\{ij\in E\}}+1+\one_{\{ij\in H\}},
\]
so \(E\) and \(H\) are complementary cycles.  This proves uniqueness up to
relabeling, without enumerating designs.  Choose
\[
 E=\{14,15,23,25,34\},\qquad
 H=\{12,13,24,35,45\}.
\]
The five memberships through block \(6\) are \(6e\), \(e\in E\), and the
other five are the complements in \(\{1,\ldots,5\}\) of the edges in \(H\).
Thus a forced normal form for the ten points is
\begin{equation}\label{eq:six-design}
 123,124,135,236,146,156,245,256,345,346.
\end{equation}
This is one normal form derived from the two cycles, not a list of candidate
supports.

Invert at the point \(123\).  The three incident circles become distinct
lines.  Their pairwise second selected intersections \(124,135,236\) are
distinct, so they form a nondegenerate triangle.  Write
\[
\begin{aligned}
 A&=124,\quad B=135,\quad C=236,\\
 X&=146,\quad Y=156,\quad U=245,\quad
 V=256,\quad W=345,\quad Z=346.
\end{aligned}
\]
Then \(X,Y\in AB\), \(U,V\in AC\), and \(W,Z\in BC\).  The three
nonincident circles have full traces
\begin{equation}\label{eq:six-circles}
 AXUWZ,\qquad BYUVW,\qquad CXYVZ.
\end{equation}

Parametrize the sides by
\[
\begin{aligned}
 AB:&\ (A,B,X,Y)=(0,1,x,y),\\
 AC:&\ (A,C,U,V)=(0,1,u,v),\\
 BC:&\ (B,C,W,Z)=(0,1,w,z)
\end{aligned}
\]
and again put \(a^2=|BC|^2\), \(b^2=|CA|^2\), \(c^2=|AB|^2\).
Power at the two omitted vertices for each circle in
\eqref{eq:six-circles} gives all six identities
\begin{align}
 c^2(1-x)&=a^2wz,&
 b^2(1-u)&=a^2(1-w)(1-z),\label{eq:six-power-a}\\
 c^2y&=b^2uv,&
 a^2(1-w)&=b^2(1-u)(1-v),\label{eq:six-power-b}\\
 c^2xy&=b^2v,&
 c^2(1-x)(1-y)&=a^2z.\label{eq:six-power-c}
\end{align}
The nine image points are distinct, so each of
\[
 x,y,u,v,w,z,\quad1-x,1-y,1-u,1-v,1-w,1-z
\]
is nonzero; every following division is legal.  Dividing the fifth identity
by the third, the sixth by the first, and combining the second with the
fourth gives
\begin{equation}\label{eq:six-substitutions}
 x=\frac1u,\qquad1-y=\frac1w,\qquad
 (1-v)(1-z)=1.
\end{equation}
The first, fourth, and third identities, with their squared side lengths
cancelled, also give the coefficient-free cycle
\begin{equation}\label{eq:six-cycle}
 \frac{(1-x)(1-w)uv}{wz(1-u)(1-v)y}=1.
\end{equation}
Substitution from \eqref{eq:six-substitutions} makes the left side of
\eqref{eq:six-cycle} equal to
\[
 \frac{((u-1)/u)(1-w)uv}
 {w\,(v/(v-1))(1-u)(1-v)((w-1)/w)}=-1.
\]
Every displayed factor is nonzero by the preceding guard.  This
contradiction proves that \(q=6\) is impossible over \(\R\).
\end{proof}

Now take \(n=10,m=5\) and write \(q=c_5\).  Circle intersection gives the
hypothesis of Lemma~\ref{lem:six-five}, so \(q\leq5\).  Equations
\eqref{eq:T} and \eqref{eq:M} become
\begin{equation}\label{eq:n10-arithmetic}
 120=c_3+4c_4+10q,\qquad3c_3-5q\geq30.
\end{equation}
Modulo four, \(c_3\equiv-2q\pmod4\).  For \(0\leq q\leq5\), every candidate
through \(8+2q\) violates the second inequality, since
\(3(8+2q)-5q=24+q<30\); all smaller integers in the same residue class do
as well.  The next congruent integer is \(12+2q\), so
\(c_3\geq12+2q\), and elimination of \(c_4\) gives
\begin{equation}\label{eq:n10-39}
 C=30+\frac34c_3-\frac32q\geq39
 \qquad(n=10,m=5).
\end{equation}

\section{The fixed small dispatch and conclusion}

For \(4\leq n\leq10\), every possible largest trace size is one of
\(m=3,\ldots,n-1\).  In that order, the lower bounds obtained from
\eqref{eq:T}, \eqref{eq:m4-values}, \eqref{eq:fano-jump},
\eqref{eq:n10-39}, and \eqref{eq:P} are
\begin{align*}
 n=4 &: \quad (4),\\
 n=5 &: \quad (10,7),\\
 n=6 &: \quad (20,11,11),\\
 n=7 &: \quad (35,17,19,16),\\
 n=8 &: \quad (56,20,25,28,22),\\
 n=9 &: \quad (84,30,29,37,40,29),\\
 n=10&: \quad (120,39,39,43,55,53,37).
\end{align*}
This is a seven-line substitution into proved formulas, not a catalogue of
configurations or supports.  Taking the minimum in each line gives
\[
 4,7,11,16,20,29,37.
\]
Construction~\ref{con:near-circle} attains every value except at eight;
Construction~\ref{con:cross} attains twenty at eight.  Together with
\eqref{eq:large-result}, this proves Theorem~\ref{thm:v3}.

\section*{Computational boundary}

No computation is a premise of this proof.  Exact-arithmetic companion
programs can corroborate the polynomial factorizations
\eqref{eq:slack}--\eqref{eq:large-odd}, the two product-fibre identities,
the forced incidence normal forms, and the final sign cancellations.  They
do not formalize projective Euler topology, inversion, directed power, or
Menelaus, and they do not establish any geometric incidence used above.
Conversely, this text uses no finite scan in \(n\) or \(m\), no generated
configuration or support family, no orbit reduction, and no solver output.

\end{document}
